

\begin{abstract}
We show how a driven-dissipative cavity coupled to a collective ensemble of atoms can dynamically generate metrologically useful spin-squeezed states. In contrast to other dissipative approaches, we do not rely on complex engineered dissipation or input states, nor do we require tuning the system to a critical point. Instead, we utilize a strong symmetry, a special type of symmetry that can occur in open quantum systems and emerges naturally in systems with collective dissipation, such as superradiance. This symmetry  preserves coherence and allows for the accumulation of an atom number-dependent Berry phase which in turn creates spin-squeezed states via emergent one-axis twisting dynamics. This work shows that it is possible to generate entanglement in an atom-cavity resonant regime with macroscopic optical excitations of the system, going beyond the typical dispersive regime with negligible optical excitations often utilized in current cavity-QED experiments. 


\pacs{}

\maketitle


An important goal in quantum metrology is to develop methods for generating quantum entanglement that reduce measurement uncertainty beyond the standard quantum limit, a fundamental bound  which defines the maximum precision of phase estimation achievable with uncorrelated particles. 
One of the most common approaches to generating metrologically-useful entanglement is via cavity-QED systems, where the cavity photons mediate highly-collective spin interactions between $N$ particles, whose resulting collective interactions are ideal for preparing spin-squeezed states \cite{Kitagawa1993, Wineland1992, Wineland1994, Pezze2018}. 

Simultaneously, extensive efforts have been devoted to the storage and protection of quantum information, including topologically-protected ground-state manifolds \cite{Kitaev2003,Nayak2008} and quantum error correction \cite{Shor1995,Terhal2015}. 
These have been extended to open quantum systems in the form of passive error correction \cite{Lidar1998,Paz1998,Barnes2000,Sarovar2005,Mirrahimi2014,Cohen2014,Moloney2015,Kapit2015,Kapit2016,Puri2017,Reiter2017,Lihm2018,Lieu2022}, which was recently shown \cite{Lieu2020, Liu2023} to be a consequence of spontaneous breaking of a ``strong symmetry'' \cite{Buca2012}, a symmetry whose conservation laws are satisfied at the single-trajectory level, which generally support multiple steady states as a result \cite{Buca2012,suppcite,Albert2014, Manzano2018,SanchezMunoz2019,Dutta2020,Minganti2023,Li2023,Zhang2024}, as opposed to ``weak'' \cite{Buca2012} or ``average'' \cite{Ma2023a, Ma2023b} symmetries which only exhibit conservation laws on average in the density matrix (see Fig.~\ref{fig:schematic}a). 
Formally, a strong symmetry is distinguished by the fact that it is an explicit symmetry of the jump operators, rather than a symmetry of the Lindblad superoperators as in the case of a weak symmetry,  ensuring the symmetry's conservation under quantum jumps.

\begin{figure}[h!]
\vspace{.12cm}
    \centering
    \includegraphics[scale=1]{DissSqFig1Half.pdf}
    \caption{
    (a) Qualitative illustration of symmetry breaking for two quantum trajectories of the wavefunction (solid, dashed). Weak: Quantum jumps move between two symmetry sectors (light, dark); symmetry breaking conserves symmetry on average by conserving the probability to be in a given sector \cite{traj}.
    Strong: Symmetry conserved for each trajectory; symmetry breaking ensures the environment cannot distinguish which sector a given trajectory corresponds to, preserving coherences between sectors.
    (b) An array of three level  atoms ($|{\uparrow}\rangle,|{\downarrow}\rangle,|e\rangle) $ are loaded in a driven  optical cavity  resonant with the  $|{\downarrow}\rangle-|e\rangle$ transition. 
    (c) Schematic of squeezing generation, see text for details. 
    \label{fig:schematic}
    }
\end{figure}

We present a protocol to dynamically generate spin squeezing in a driven-dissipative optical cavity via a Berry phase using a strong symmetry. 
We consider a spin-1/2 system composed of long-lived states $|{\uparrow}\rangle,|{\downarrow}\rangle$ (e.g., clock/ground states). In its simplest implementation, the system is initialized in an equal superposition of $|{\uparrow}\rangle,|{\downarrow}\rangle$ and an external laser is used to drive $|{\downarrow}\rangle$ to an excited state $|e\rangle$, which is resonant with the cavity (see Fig.~\ref{fig:schematic}b).  Here, the strong symmetry corresponds to global phase rotations on the $|{\downarrow}\rangle, |e\rangle$ states, and the combined number of atoms in the two states is conserved, which emerges because $|{\uparrow}\rangle$ is not coupled to the other states.
The protocol consists of three steps (see Fig.~\ref{fig:schematic}c). (i) A resonant drive is turned on, and the interplay between coherent drive and collective dissipation brings the system to a steady state which defines an effective qubit state $|1\rangle$, composed of a superposition of $|{\downarrow}\rangle, |e\rangle$, and $|0\rangle \equiv |{\uparrow}\rangle$.
(ii) The drive is then detuned from resonance so the phase of the drive changes in time, leading to a corresponding change in the phase of the steady state, simultaneously resulting in the accumulation of a global phase on $|{\downarrow} \rangle, |e\rangle$ (which form $|1\rangle$) relative to $|{\uparrow}\rangle$.
(iii) After the drive is turned off, the accumulated phase is preserved as $|e\rangle$ relaxes to $|{\downarrow}\rangle$, i.e., $|1\rangle \to |{\downarrow}\rangle$, due to the strong symmetry preserving relative phases between ($|{\downarrow}\rangle, |e\rangle$) and $|{\uparrow}\rangle$.

The accumulated phase on $|{\downarrow}\rangle$ is thus a dissipative Berry phase $\phi_B$ defined by the closed loop formed by the steady-state Bloch vector. We utilize a gauge where the accumulation of $\phi_B$ only occurs during (ii), so we focus primarily on the dynamics of this step unless otherwise noted. 
Because $\phi_B$ is sensitive to variations in the number of initial $|{\downarrow}\rangle$ atoms, the Berry phase generates effective one-axis twisting (OAT) dynamics, a paradigmatic fully connected Ising model ($\hat S_z^2$) \cite{Kitagawa1993}, 
which shears an initial coherent spin state distribution into an entangled state with scalable (in $N$) spin squeezing.
The strong symmetry helps preserve the accumulated phase throughout the dynamics and allows for the turn on/off of the drive to be done quickly \cite{suppcite}. Steps (i,ii) can be combined due to the rapid relaxation to the steady state $N \Gamma \gg \delta$. 

We show that our protocol achieves scaling comparable to or even better than state-of-the-art spin squeezing generation protocols in the presence of realistic sources of decoherence. 
Given that our protocol is not limited to the weak-excitation  regime of many coherent \cite{Schleier-Smith2010, Borregaard2017, Pedrozo-Penafiel2020, Colombo2022, Borregaard2017a, Luo2024, Leroux2010, Braverman2019, Barberena2023b}, measurement-based \cite{Schleier-Smith2010a, Cox2016, Hosten2016, Thomsen2002, Chen2014,Bao2020, Bowden2020, Robinson2022, Greve2022, Barberena2023b}, and dissipative \cite{DallaTorre2013} approaches, it enables much faster squeezing generation timescales, which can reduce the detrimental effect of other types of decoherence processes beyond those intrinsic to the cavity setup. Unlike most dissipative schemes which achieve spin squeezing in the steady state of a driven cavity near criticality \cite{DallaTorre2013, Lee2014a, Fernandez-Lorenzo2017, Groszkowski2022,Pavlov2023, Barberena2024,Gonzalez-Tudela2013, Wolfe2014, Somech2024}, our protocol dynamically generates squeezing away from criticality and stored in metrologically-relevant states (e.g., a clock transition) after the drive is turned off. Finally, unlike measurement-based approaches \cite{Thomsen2002,Schleier-Smith2010a, Chen2014, Cox2016, Hosten2016, Bao2020, Bowden2020, Robinson2022, Greve2022, Barberena2023b}, our protocol is deterministic and not limited by detection efficiency.

\emph{Model}.---We consider a system composed of a single driven lossy cavity mode $\hat a$ that resonantly interacts with $|{\downarrow} \rangle$ and an excited state $|e \rangle$ with single-photon Rabi frequency $g_c$ and power decay rate $\kappa$ (see Fig.~\ref{fig:schematic}b). We focus on the bad cavity limit $\kappa \gg \sqrt{N} g_c$,
so we can adiabatically eliminate the cavity \cite{suppcite}.
The spin dynamics of the system are described by the master equation
\begin{subequations}
\begin{equation}
\label{eq:spin_model}
\partial_t \hat\rho = - i [\hat H_{B,\delta} ,\hat\rho] + \Gamma \left( \hat l \hat\rho \hat l^\dagger - \frac{1}{2} \{\hat l^\dagger \hat l, \hat\rho \} \right),
\end{equation}
\begin{equation}
    \hat H_B = \frac{\Omega \hat J^+ + \Omega^* \hat J^-}{2}, \quad \hat H_\delta = |\Omega| \hat J_x - \delta \hat N_e,
\end{equation}
\end{subequations}
for an effective Rabi frequency $\Omega$ due to the cavity drive. The Hamiltonians correspond to the rotating frame of the atomic $|{\downarrow}\rangle \to |e\rangle$ transition frequency ($\hat H_B$) or the drive frequency ($\hat H_\delta$), where $\hat J^- \equiv \sum_i |{\downarrow}_i\rangle \langle e_i|$ defines an effective two-level collective operator in the relevant frame, $\hat N_e \equiv \sum_i |e_i\rangle \langle e_i|$, and $\Gamma=4 g_c^2/\kappa$ is the collective decay rate with Lindblad jump operator $\hat l = \hat J^-$. 

The two equivalent versions of the Hamiltonian support two different perspectives for the dynamics. $\hat H_B$ supports a geometrical picture where slow  variations in $\Omega$, via
$\Omega = |\Omega| e^{i \delta t}$,
result in the accumulation of a Berry phase. Alternatively, the time-dependence is eliminated in the rotating frame of the drive for $\hat H_\delta$, where $\delta$ physically corresponds to the detuning of the drive (see Fig.~\ref{fig:schematic}b). The latter further helps to clarify non-adiabatic effects in the dynamics, including decoherence. Throughout, we will use $\hat H_\delta$ for its convenience for calculations.

\emph{Mean field theory}.---The two-level piece $(|{\downarrow}\rangle-|e\rangle)$ of the above model has previously been studied in several contexts \cite{Drummond1978, Puri1980, Walls1980, Kilin1980, Iemini2018, Hannukainen2018, Barberena2019, Link2019, Somech2024, Barberena2024,Gonzalez-Tudela2013, Wolfe2014}, so we summarize key features assuming $N_J$ total atoms within this manifold. First, we recast the spin operators in terms of Schwinger bosons $\hat d, \hat e$, bosonic operators that annihilate atoms in the $|{\downarrow}\rangle$ and $|e\rangle$
states. In terms of these bosonic operators,  $\hat{J}^- \equiv \hat d^\dagger \hat e ,\hat{J}_z \equiv (\hat e^\dagger \hat e - \hat d^\dagger \hat d)/2$. In the large $N_J$ limit, the mean-field equations are
\begin{subequations}
    \begin{gather}
        \partial_t d = - i \frac{\Omega}{2} e + i \omega_B[N_J] d + \frac{\Gamma}{2} |e|^2 d,\\
        \partial_t e = - i \frac{\Omega}{2} d + i (\delta + \omega_B[N_J ]) e - \frac{\Gamma}{2} |d|^2 e,
    \end{gather}
\end{subequations}
where $d \equiv e^{i \omega_B[N_J] t} \langle \hat d\rangle, e \equiv e^{i \omega_B[N_J] t} \langle \hat e \rangle$. Here, we have gone to a rotating frame defined by $\omega_B[N_J]$, which is a dynamical phase that emerges due to Berry phase accumulation. It is determined by requiring that the above differential equations reach a steady state in the rotating frame \cite{suppcite}
\begin{equation}
    \omega_B[N_J ] = \frac{\delta}{2 \cos \theta_J[N_J]} - \frac{\delta}{2},
\end{equation}
where $\cos \theta_J[N_J] \equiv - 2\langle \hat J_z \rangle/N_J$. Physically, this describes the accumulation of a single-particle global phase on $|{\downarrow}\rangle, |e\rangle$.

When $\delta = 0$, the two-level system features spin-polarized and mixed steady-state phases separated by a  critical point  at $\Omega_c[N_J] = \frac{N_J \Gamma}{2}$, which was recently observed for the first time \cite{Song2024}.
When $\Omega > \Omega_c$, the system is in a mixed phase which exhibits zero steady state inversion ($\langle \hat J_z \rangle = 0$). In this phase, the drive overwhelms the collective dissipation, leading to modified Rabi cycles whose competition with the collective dissipation causes the spins to decohere. A Berry phase cannot be defined in such conditions.

When $\Omega < \Omega_c$, the system is  in the spin-polarized phase. The steady state exhibits nonzero inversion and collective coherence due to the competition of the drive and dissipation, so there are no photons inside the cavity in this phase, $\langle \hat a \rangle = i \frac{g_c}{\kappa/2} (\langle \hat J^- \rangle + i \Omega/\Gamma) = 0$ \cite{suppcite}. For large $N_J$, the steady state  is  highly  pure and admits an excellent  mean-field description:
\begin{equation}
     |\psi_{N_J}\rangle_\text{ss}  \approx 
     \left(\cos \frac{\theta_J[N_J]}{2} |{\downarrow}\rangle + e^{-i \phi_J[N_J]} \sin \frac{\theta_J[N_J]}{2} |e \rangle\right)^{N_J},
\end{equation}
where $\cos \theta_J[N_J] = \sqrt{1 - |\Omega/\Omega_c[N_J]|^2}$, $\phi_J[N_J] = \pi/2$. When $0 < |\delta| \ll N_J \Gamma$, $\phi_J[N_J]$ is slightly perturbed, leading to small photon buildup in the cavity.  In this case, changes in the phase of $\Omega$ induce equivalent changes in $\phi_J$, corresponding to step (ii) in the protocol. The steady state in this regime is illustrated in Fig.~\ref{fig:BlochSpheres}a.

\begin{figure}
    \centering
    \includegraphics{FigBloch.pdf}
    \caption{Bloch sphere description. (a) Bloch sphere dynamics for two-level model $\hat J$. The Berry phase accumulation process is shown for each step of the protocol. (b) Resulting dynamics on effective $\hat S$ Bloch sphere formed from the $|0\rangle, |1\rangle$ states during step (ii).}
    \label{fig:BlochSpheres}
\end{figure}

\emph{Strong symmetry}.---Now consider an initial superposition $|\psi(0)\rangle = ( |{\uparrow}\rangle/\sqrt{2} + e^{i \varphi} |{\downarrow}\rangle/\sqrt{2})^N$. Since the $|{\uparrow}\rangle$ state is not coupled to the $\{|{\downarrow}\rangle, |e\rangle \}$ states, the state in the spin-polarized phase during (ii) becomes
\begin{equation}
\label{eq:steadypsi}
 |\psi(t) \rangle \approx \sum_{N_J = 0}^N \sqrt{\binom{N}{N_J}} \frac{ e^{i N_J \varphi}}{2^N}  e^{i \phi_B[N_J,t]} \mathcal{S}|\psi_{N_J} \rangle_\text{ss} |{\uparrow}\rangle^{N_{\uparrow}}, 
\end{equation} 
where $N_{\uparrow} = N - N_J$, $\mathcal{S}$ symmetrizes the $\{|{\downarrow}\rangle, |e\rangle \}$ and $|{\uparrow}\rangle$ states, and $\phi_B[N_J,t]$ captures the $N_J$ and time dependence of the Berry phase accumulation, which reduces to $\omega_B t$ at the single-particle mean-field level. 

To understand the dynamics, we define a basis $|0\rangle \equiv |{\uparrow} \rangle, |1 \rangle \equiv \cos \frac{\tilde{\theta}_J}{2} |{\downarrow} \rangle + e^{-i \tilde{\phi}_J} \sin \frac{\tilde{\theta}_J}{2} |e\rangle$ with collective spin $\hat S$, where tildes denote mean-field values at $N_J = \langle \hat N_J \rangle = N/2$. In this basis, $ \mathcal{S}|\psi_{N_J} \rangle |{\uparrow}\rangle^{N_{\uparrow}}$ describe collective Dicke states of $\hat S$ with $m_S = \frac{N_{\uparrow} - N_J}{2}$, so $\hat S_z \approx \frac{N}{2} - \hat N_J$, and the evolution of $\phi_B[N_J,t]$ approximately corresponds to phase evolution for a specific $m_S$ state. Since $N_J$ for each term in the superposition is conserved, they
evolve independently, accumulating different Berry phases. This in turn leads to a shearing of the noise distribution as observed in standard OAT dynamics, which similarly preserves the total magnetization (see Fig.~\ref{fig:BlochSpheres}b).

The $\hat S$ picture is possible because in the spin-polarized phase, a ``strong'' symmetry \cite{Buca2012} undergoes spontaneous symmetry breaking.
Strong symmetries leave the Lindblad jump operators invariant, enforcing conserved quantities on an individual trajectory level \cite{Buca2012, SanchezMunoz2019, Lieu2020, Dutta2020,suppcite,Albert2014, Manzano2018,Zhang2024, Li2023,Liu2023,Ma2023a, Ma2023b, Minganti2023}. Here, the symmetry corresponds to phase rotations $|{\downarrow}_i\rangle, |e_i\rangle \to e^{i \varphi'} |{\downarrow}_i\rangle, e^{i \varphi'} |e_{i}\rangle$ for all $i$ and number conservation $\hat N_J \equiv \sum_i |{\downarrow}_{i} \rangle \langle {\downarrow}_{i}| + |e_{i} \rangle \langle e_{i}|$, where eigenvalues of $\hat N_J$ define the symmetry sectors. Conservation of $\hat N_J$ implies conservation of the $|{\uparrow}\rangle$ population, while strong-symmetry breaking preserves the coherences between terms with different $N_J$ \cite{Lieu2020,Liu2023}.

Physically, this means for any quantum trajectory in one sector (e.g., states with exactly $N_J$ atoms in $|{\downarrow}\rangle, |e\rangle$), equivalent trajectories exist for the other sectors ($N_J' \neq N_J$ atoms in $|{\downarrow}\rangle, |e\rangle$) and thus are indistinguishable from the perspective of the environment. For a weak symmetry, quantum jumps could take $N_J$ states to $N_J'$ states and vice-versa, and $\hat N_J$ could only be conserved on average (see Fig.~\ref{fig:schematic}a). Thus, the strong symmetry prevents the environment from projecting the system into a specific sector, thereby preserving the coherences between sectors. Formally, the steady-state space forms a decoherence-free subspace \cite{Lidar1998}, although strong symmetries can also result in more general noiseless subsystems \cite{Knill2000, Albert2014, Lieu2020}.

The  strong symmetry further ensures that the squeezing generated outside the $|{\uparrow}\rangle ,|{\downarrow} \rangle$ manifold is mapped back to $|{\uparrow}\rangle,|{\downarrow}\rangle$ by simply turning off the drive. 
All that occurs is $\mathcal{S}|\psi_{N_J} \rangle |{\uparrow}\rangle^{N_{\uparrow}} + \mathcal{S}|\psi_{N_J'} \rangle |{\uparrow}\rangle^{N_{\uparrow}'} \to \mathcal{S}|{\downarrow} \rangle^{N_J}|{\uparrow}\rangle^{N_{\uparrow}}+\mathcal{S}|{\downarrow} \rangle^{N_J'}|{\uparrow}\rangle^{N_{\uparrow}'}$ with no state mixing when $|N_J - N_J'| < \mathcal{O}(\sqrt{N})$  \cite{suppcite}, so the Berry phase evolution is retained. 
This allows us to consider dynamics with significant $|e\rangle$ excitations while nevertheless generating squeezing in the desired $|{\downarrow}\rangle, |{\uparrow}\rangle$ manifold.

{\it Squeezing dynamics}.---We are interested in deriving an effective Hamiltonian from the Berry phase $\hat {H}_{\text{eff}} |\Psi_{N_J} \rangle = (-\partial_t  \phi_B) |\Psi_{N_J}\rangle$, where $|\Psi_{N_J} \rangle \equiv \mathcal{S}|\psi_{N_J} \rangle_\text{ss} |{\uparrow}\rangle^{N_{\uparrow}}$.
We note that, for large $N_J$,  the  accumulated  Berry phase is that of a large spin in a magnetic field and  given by the subtended solid angle, $\frac{N_J}{2}(1- \cos \theta_J)$. Therefore,   $\partial_t \phi_B =  -\delta [N_J (1-\sqrt{1-\sin^2 \theta_J})]/2$. Since for small $\delta$, $\sin \theta_J \approx  \Omega/(N_J \Gamma/2)$, we incorporate quantum fluctuations in $\phi_B$ by replacing  $N_J \to \hat{N}_J = N/2 - \hat S_z$  and  $\sin \theta_J \to \sin \theta_J[\hat{N}_J]\approx \frac{N/2}{\hat N_J} \sin \tilde{\theta}_J$ in $\partial_t \phi_B$. 
Expanding in powers of $|\hat S_z/N|\ll 1$, $\hat {H}_{\text{eff}}$ becomes 
\begin{equation}
\label{eq:Heff}
     \hat H_{\text{eff}} \approx  -\tilde{\omega}_B \hat S_z + \check{\chi} \hat S_z^2 + \cdots, \quad \check{\chi} = - \frac{\delta \sin^2 \tilde{\theta}_J}{2 N \cos^3 \tilde{\theta}_J},
\end{equation} where we have dropped a constant term as a global phase. Here, we see that variations of the Berry phase with $N_J$ give rise to OAT dynamics and describe the mean-field rotating frame of $\tilde{\omega}_B$.

When $\delta \neq 0$, the cavity's steady state is no longer the vacuum. The finite  coherence of the cavity,  proportional to the modified jump operator $\hat l = \hat J^- + i \Omega/\Gamma \approx -i (\kappa/2 g_c) \langle \hat a \rangle$, thus   depends on  the population of the  various $N_J$ terms, which radiate at different rates,   leading to dephasing. 
By using its  corresponding  mean-field value at linear order in $\delta$, $ \langle \hat l \rangle = N_J \left(\frac{\delta}{2 \Omega} \sin \theta_J\tan \theta_J - i \frac{N_J \Gamma}{4 \Omega} \sin^2 \theta_J \right) + i \Omega/\Gamma $, and  replacing  $N_J$ and $\sin\theta_J$ with their operator values as above, we obtain  
\begin{equation}
\label{eq:l}
    \hat l\approx \frac{\delta}{\Gamma} \tan \tilde{\theta}_J + \frac{2 \delta  \sin \tilde{\theta}_J}{N \Gamma \cos^3 \tilde{\theta}_J } \hat S_z + \cdots .
\end{equation} 
The linear $\hat S_z$ term indicates the presence of an effective collective dephasing due to the environment's ability to slowly gain information about $\hat N_J$. The constant term does not affect the $\hat S$ dynamics. The ratio of the OAT and collective dephasing rates is $N \Gamma \cos^3 \tilde{\theta}_J/(8 \delta)$, so the effect of collective dephasing is reduced when the detuning is small, corresponding to the adiabatic limit where everything is determined by the Berry phase accumulation and there is minimal distinguishability between the different $N_J$ terms in the state. This scaling is consistent with recent results investigating the manipulation of these steady-state spaces \cite{Santos2024}. Further details for both derivations are presented in the End Matter.

To extend the analysis beyond the small $\delta$ limit, we utilize a generalized Holstein-Primakoff approximation, from which we find \cite{suppcite}
\begin{subequations}
\label{eq:AEHP}
    \begin{equation}
        \hat H_{\text{eff}} =  -\frac{\delta}{2N } \frac{N^2 \Gamma^2 \sin \tilde{\theta}_J \tan \tilde{\theta}_J}{N^2 \Gamma^2 \cos^2 \tilde{\theta}_J + 16 \delta^2 \sec^2 \tilde{\theta}_J} \hat S_z^2,
    \end{equation}
    \begin{equation}
        \hat l \approx 2  e^{-i \tilde{\phi}_J} \frac{\delta \tan \tilde{\theta}_J (i N \Gamma - 4 \delta \sec \tilde{\theta}_J ) }{N^2 \Gamma^2 \cos^2 \tilde{\theta}_J + 16 \delta^2 \sec^2 \tilde{\theta}_J} \hat S_z,
    \end{equation}
\end{subequations}
where the constant term has been dropped from the jump operator. We benchmark these equations with quantum trajectories in the End Matter.

Most typical squeezing protocols are restricted to the perturbative, weak-excitation regime where $|e\rangle$ is not strongly populated and can be adiabatically eliminated \cite{Schleier-Smith2010, Borregaard2017, Pedrozo-Penafiel2020, Colombo2022, Borregaard2017a, Luo2024, Leroux2010, Braverman2019, Barberena2023b,Mivehvar2021}. However, the dynamics become correspondingly slower due to the reduced effective atom-cavity interaction. The slower dynamics  in the weak-excitation limit can be seen  in our protocol via the fact that both the OAT and dephasing rates are proportional to $\sin^2 \tilde{\theta}_J$, and  thus they become slow for weak excitations since  $\sin^2 \tilde{\theta}_J \to 0$. Nevertheless, our protocol extends beyond this regime, allowing for a  significant acceleration of the atom-cavity  dynamics. 
Therefore, while the scaling of our protocol with $N$ is the same as other OAT approaches, our protocol enables a dramatic reduction in the timescales necessary to achieve the same amount of squeezing by going beyond the perturbative regime.

{\it Single particle decoherence}.---Depending on the choice of states, there are different types of effective single particle decoherence in $\hat{S}$. A fundamental source of decoherence is spontaneous emission, which enters the master equation via
$\gamma_{e\updownarrow} \sum_i \left(\hat \sigma_i^{\updownarrow e} \hat \rho \hat \sigma_i^{e \updownarrow} - \frac{1}{2} \{\hat \rho , \hat \sigma_i^{ee}\} \right)$, where $ \hat \sigma^{ab} \equiv |a \rangle \langle b|$,
describing single particle  decay from $|e \rangle$ to either $|{\updownarrow}\rangle =|{\uparrow}\rangle$ or $ |{\downarrow} \rangle$, each of which has different effects on the effective $\hat S$ Bloch sphere dynamics. We model spontaneous emission to $|{\uparrow}\rangle$ at rate $\gamma_{e {\uparrow}}$ as an effective spin flip from $|1\rangle$ to $|0 \rangle$ at rate $\gamma_- \equiv \gamma_{e{\uparrow}} \sin^2 \frac{\tilde{\theta}_J}{2}$, which is reduced from $\gamma_{e{\uparrow}}$ by the excited state fraction of $|1 \rangle$. 

 
We treat spontaneous emission to $|{\downarrow} \rangle$ at rate $\gamma_{e \downarrow}$ as an effective dephasing process in the reduced $|0\rangle, |1\rangle$ basis, which enters the effective master equation for $\hat S$ as $\gamma_d \sum_i \left(\hat \sigma_i^{11} \hat \rho_S \hat \sigma_i^{11} - \frac{1}{2} \{\hat \rho_S , \hat \sigma_i^{11}\} \right)$
with $\gamma_d = \gamma_{e \downarrow} \sin^2 \frac{\tilde{\theta}_J}{2}$. However, given that in typical experiments there are other sources of effective single-particle dephasing, in the discussion below we  assume a $\tilde{\theta}_J$-independent rate. If $|{\downarrow}\rangle \equiv {^1S_0}, |e\rangle \equiv {^3 P_1}, |{\uparrow} \rangle \equiv {^3 P_0}$ for $^{87,88}$Sr, creating squeezing on a clock transition, then $\gamma_{e{\uparrow}} = 0$, and there are no effective spin flips. In contrast, effective single-particle dephasing is always present.

We quantify the squeezing via the Wineland squeezing parameter $\xi$ \cite{Wineland1992, Wineland1994}, which describes the reduction in phase measurement uncertainty beyond the standard quantum limit and is defined as $\xi^2 = \frac{N \text{min} (\Delta \hat{S}_{\perp})^2}{|\langle \hat{\mathbf{S}} \rangle|^2}$, where $\text{min}(\Delta \hat{S}_{\perp})^2$ denotes the minimum variance in directions perpendicular to $\langle \hat{\mathbf{S}} \rangle$. We analytically optimize the squeezing via an expansion of the $\hat S$ dynamics \cite{suppcite, Chu2021, Barberena2023b}
\begin{equation}
    \xi^2(t) \approx e^{\gamma_d t}\left(\frac{e^{\gamma_d t} + N \check{\Gamma} t}{N^2 \check{\chi}^2 t^2} + \frac{N^2 \check{\chi}^4 t^4}{6} + \frac{2}{3} \gamma_- t\right),
\end{equation}
where $\check \chi, \check \Gamma$ denote the OAT and collective dephasing rates, respectively. 
In Table \ref{tab:optsq}, we optimize the squeezing when one source of single-particle dephasing dominates.
As an example, in Fig.~\ref{fig:AEHP} we plot $\xi^2_{\text{opt}}$ and $t_{\text{opt}}$
for $N=10^6$, with either only effective spin flips or only effective single-particle dephasing, determined via an exact solution \cite{Chu2021} using Eq.~(\ref{eq:AEHP}) for the collective dynamics. 

\begin{table}
\centering
\def\arraystretch{1}
\setlength\tabcolsep{3.5mm}
\begin{tabular}{lcc}
\toprule
\toprule
   & $\xi^2_{\text{opt}}$ & $t_{\text{opt}}$ \\
\midrule
 Dephasing & $\propto \frac{1}{N^{2/3}}$ & $\min\left[\frac{3.04}{\sqrt{\Gamma /\gamma_d} N^{1/6}},1\right] \times \gamma_d^{-1}$\\
 Spin flips &  $ \frac{4\sqrt{2/3}}{\sqrt{N \Gamma/\gamma_{e{\uparrow}}}}$ & $\frac{\sqrt{6}}{ \sin^2 \frac{\tilde{\theta}_J}{2} \sqrt{N \Gamma/\gamma_{e{\uparrow}}}} \times \gamma_{e \uparrow}^{-1}$ \\
 \bottomrule
 \bottomrule
\end{tabular}
\caption{Optimal performance of squeezing protocol in the presence of single-particle dephasing. Similar scaling is realized by weak-excitation OAT protocols ($\sin^2 \frac{\tilde{\theta}_J}{2} \ll 1$) up to constant factors \cite{Schleier-Smith2010, Borregaard2017, Lewis-Swan2018, suppcite}. Scaling comparisons for other protocols and experimental values are reported in the supplement \cite{suppcite}. \label{tab:optsq}}
\end{table}

For spin flips ($\gamma_- = \Gamma \sin^2\frac{\tilde{\theta}_J}{2}, \gamma_d = 0$), the optimal squeezing is $24.4$ dB at $\delta = 0.014 N \Gamma$ in the limit $\sin^2\tilde{\theta}_J/2 \to 0$, i.e., $\Omega \to 0$ (diamond). However, the squeezing time scales grow like $1/\sin^{2} \frac{\tilde{\theta}_J}{2}$, so additional single-particle dephasing always shifts the optimal drive to stronger $\Omega$. 
Increasing the drive to $\cos \tilde{\theta}_J = 0.7 $, the squeezing at the optimal $\delta = 0.005 N \Gamma$ is $24.0$ dB in time $0.0163/\Gamma$ (square), which is 28 times faster than at the optimal detuning for $\cos \tilde{\theta}_J = 0.99$ (weak-excitation regime), so additional single-particle dephasing can be mitigated with minimal reduction in the squeezing.

For single-particle dephasing ($\gamma_- = 0, \gamma_d = 100 \Gamma$), the optimal squeezing is 23.1 dB in time $0.0085/\Gamma = 0.85/\gamma_d$, which occurs in the limit $\delta \to 0$ with $\cos^3 \tilde \theta_J \approx 33 \delta/N \Gamma$ (circle). However, in this limit finite-size effects which reduce the validity of Eq.~(\ref{eq:AEHP}) occur, so this limit cannot be realized. Nevertheless, moving away from criticality to $\cos \tilde{\theta}_J = 0.6$, which has optimal $\delta = 0.01 N \Gamma$ (triangle), there is again minimal reduction in the squeezing (21.3 dB) and minimal increase in the squeezing time ($0.009/\Gamma = 0.9/\gamma_d$).

When both types of decoherence are relevant, optimal squeezing occurs at intermediate drives. For $\gamma_- = \Gamma \sin^2\frac{\tilde{\theta}_J}{2}, \gamma_d = 100 \Gamma$, the optimal squeezing occurs at $\delta = 0.0004 N \Gamma, \cos \tilde{\theta}_J = 0.23$ with $20.8$ dB squeezing in time $0.0044/\Gamma = 0.44/\gamma_d$ (star). 


\begin{figure}
    \centering
    \includegraphics[scale=.29]{Fig3SymbolResubShort.pdf}
    \caption{Optimal squeezing (top) and squeezing time (bottom) as a function of $\delta/N \Gamma, \cos \tilde{\theta}_J$ for $N=10^6$ using Eq.~(\ref{eq:AEHP}). (a,c) Only effective spin flips in $\hat{S}$: $\gamma_{-} = \Gamma \sin^2\frac{\tilde{\theta}_J}{2}, \gamma_d = 0$. (b,d) Only single-particle dephasing: $ \gamma_- = 0, \gamma_d = 100 \Gamma$. Black lines indicate the optimal $\delta$ for the most squeezing as a function of $\cos \tilde{\theta}_J$ using perturbative expansions. Shapes correspond to parameters discussed in main text.
    }
    \label{fig:AEHP}
\end{figure}

\emph{Outlook}.---While we have focused on generating OAT dynamics, an important question is whether other types of dynamics can be generated, such as quantum non-demolition squeezing \cite{Schleier-Smith2010a, Hosten2016, Cox2016, Barberena2023b}.
Furthermore, an interesting direction is the generalization of our protocol to multilevel systems featuring a rich dark state structure \cite{PineiroOrioli2022, Lin2022, Fan2023}, which readily realize generalized forms of the strong symmetry considered here. Finally, investigations on  how predicted behaviors are affected by realistic experimental conditions  outside the  bad cavity limit or with inhomogeneous atom-light coupling are important for its implementation in optical clocks \cite{Robinson2022}.


\begin{acknowledgments}
    We thank B.~Sundar, D.~Barberena, S.~Lieu, and J.~Hur for helpful discussions and feedback. 
 This work is supported by the VBFF,  AFOSR grants FA9550-24-1-0179,  by the NSF JILA-PFC PHY-2317149, QLCI-OMA-2016244, by the U.S. Department of Energy, Office of Science, National Quantum Information Science Research Centers Quantum Systems Accelerator, and by NIST. J.T.Y.~was supported in part by the NWO Talent Programme (project number VI.Veni.222.312), which is (partly) financed by the Dutch Research Council (NWO).
\end{acknowledgments}

\bibliography{StrongSqueezing,Extra}
